\scanpage{104}\mspage{- 115 -}
% archive sheet number: 103
toute sous-algèbre de type fini de $C\handwritten{[}T\handwritten{]}$ vérifie la condition des chaînes, il en est de même de toute sous-algèbre de type fini de $C_{\underline{r}}\handwritten{[}T\handwritten{]}$, ce qui achève la démonstration.

\uline{Corollaire 1 .- Soient} $Y$ \uline{un préschéma irréductible localement noethérien}, $X$ \uline{un préschéma irréductible}, $f:X\handwritten{\to}Y$ \uline{un morphisme} \deleted{\unsure{de type}} \uline{dominant de type fini. Soit} $\handwritten{\xi}$ (resp. $\handwritten{\eta}$) \uline{le point générique de} $X$ (resp. $Y$) \deleted{\illegible{1}}, \uline{et soit} $e=\dim(f^{-1}(\handwritten{\eta}))=\operatorname{deg.tr.}_{\handwritten{\underline{k}}(\handwritten{\eta})}\handwritten{\underline{k}}(\handwritten{\xi})$ (``dimension de la fibre générique\unsure{/}''). \uline{Pour tout} $x\handwritten{\in}X$, \uline{on a, en posant} $y=f(x)$
\[
 e+\operatorname{codim}_Y\handwritten{\overline{\{y\}}}\handwritten{\geq}
 \operatorname{deg.tr.}_{\handwritten{\underline{k}}(y)}\handwritten{\underline{k}}(x)
 +\operatorname{codim}_X\handwritten{\overline{\{x\}}}.\tag*{(13)}
\]
% slanted
\margin{\handwritten{$\dim\underline O_x\leq\dim\underline O_y+\dim\underline O_x\otimes k(y)-\delta(x)$\\
$\delta(x)=\dim_{\deleted{\unsure{y}}x} f^{-1}f(x)-e$}}
% slanted
\margin{\handwritten{formule hybride, donner plutôt \uline{deux} formules, l'une en termes d'alg. comm., l'autre en termes ``géom.''}}

\uline{En outre, les deux membres sont égaux lorsque} $\underline O_y$ \uline{est un quotient d'anneau local régulier}.
% slanted
\margin{\handwritten{\deleted{\illegible{1}} cond. des ch. \unsure{stricte}}}
\margin{\handwritten{\uuline{Variante} \uuline{au cor} $\left\{\begin{aligned}
 \dim\underline O_x&\leq\dim\underline O_y+\dim\underline O_{\deleted{\unsure{y}}x}\otimes\underline k(y)\\
 &\quad-[\dim_x f^{-1}f(x)-e]
\end{aligned}\right.$}}

Comme la question est locale, on peut supposer $X$ et $Y$ affines\inserted{\handwritten{ et réduits}}, \deleted{\illegible{11} \unsure{l'hypothèse} \illegible{5}} \inserted{soient} $X=\operatorname{Spec}(B)$, $Y=\operatorname{Spec}(A)$, où $B$ est une algèbre intègre de type fini sur $A$; comme $f$ est dominant, $A$ s'identifie à une sous-algèbre de $B$ (§ 1, nº2, cor.4 de la prop.5), donc $\underline O_y$ à une sous-algèbre de $\underline O_x$. Comme $\handwritten{\underline k(\eta)}$ et $\handwritten{\underline k(\xi)}$ sont les corps de fractions de $A$ et $B$ respectivement, donc de \deleted{$\unsure{\underline O_x}$ et} $\underline O_y$ et $\underline O_x$ respectivement, la formule (13) résulte de (12) appliquée aux anneaux $\underline O_y$ et $\underline O_x$.
% slanted
\margin{\handwritten{\deleted{Si $\bar x\subset X$ fermé irréductible, soit $\dim\unsure{Z}/f(x)$} \unsure{rédaction voir ? de suite}}}

On notera d'ailleurs que $\operatorname{deg.tr.}_{\handwritten{\underline k}(y)}\handwritten{\underline k}(x)$ n'est autre que la dimension de $f^{-1}(y)\handwritten{\cap\overline{\{x\}}}$, comme il résulte de la prop.7 du nº2, appliquée à $f^{-1}(y)$ considéré comme préschéma algébrique sur $\handwritten{\underline k}(y)$. En particulier:
\margin{\handwritten{\deleted{\illegible{1}}}}

\uline{Corollaire 2} .- \deleted{\unsure{Sous}} \uline{les hypothèses du cor.1 étant vérifiées, supposons en outre que} $x$ \uline{soit fermé dans} $f^{-1}(y)$ \uline{et que} $\underline O_y$ \uline{soit un quotient d'anneau régulier. Alors}
% slanted
\margin{\handwritten{\deleted{\illegible{1}} cond. des chaînes stricte}}
\[
 \operatorname{codim}_X\handwritten{\overline{\{x\}}}
 =e+\operatorname{codim}_Y\handwritten{\overline{\{y\}}}.\tag*{(14)}
\]

\uline{Corollaire 3 .- Les hypothèses du cor.1 étant vérifiées, on a}
